\chapter{Eosinophilic Fasciitis (Shulman Syndrome)}
\index{Eosinophilic Fasciitis (Shulman Syndrome)}
\label{sec:Eosinophilic Fasciitis (Shulman Syndrome)}

\section{Overview}
Eosinophilic fasciitis is a rare, fibrosing disorder of the fascia characterized by symmetrical, painful swelling and induration of the skin and subcutaneous tissue, with peripheral eosinophilia, with an incidence of approximately 1 per million, equal sex distribution, and typically affecting adults 40--60 years of age. This genetic disorder results from specific gene mutations or chromosomal abnormalities. It may be present at birth or manifest later in life depending on the underlying mechanism. Genetic disorders result from abnormalities in an individual's DNA, ranging from single-gene mutations to chromosomal abnormalities. These conditions may be inherited from parents or arise from new mutations. Pattern of inheritance depends on whether the mutation is dominant, recessive, or X-linked. Genetic counseling helps families understand risks and make informed decisions.
\section{Causes}
The cause is unknown, but triggers include vigorous exercise or trauma (in up to 50\%), infections, and medications (statins, antitubercular drugs), with Disease mechanism involving inflammation and thickening of the deep fascia, infiltration by lymphocytes, eosinophils, macrophages, and plasma cells, leading to collagen deposition and scarring. Genetic disorders result from mutations in specific genes that encode proteins essential for normal cellular function. The pattern of inheritance depends on whether the mutation is dominant, recessive, or X-linked.   
\section{Risk Factors}

\begin{itemize}[noitemsep]
  \item Genetic predisposition and family history
  \item Advancing age
  \item Lifestyle factors (diet, exercise, smoking, alcohol)
  \item Environmental exposures
  \item Underlying medical conditions
  \item Medication use
\end{itemize}
\section{Signs and Symptoms}

\subsection{Common symptoms}
\begin{itemize}[noitemsep]
  \item \item You may experience acute or subacute onset of symmetrical, painful swelling and induration of the arms and legs (sparing the hands, feet, and face---a distinguishing feature from systemic sclerosis), with a characteristic ``peau d'orange'' or cobblestone appearance, and the groove sign (linear depression overlying superficial veins) being a pathognomonic finding, with joint contractures developing due to fascial scarring.

  \item Pain or discomfort in the affected area
  \end{itemize}

\subsection{Emergency warning signs}
\begin{itemize}[noitemsep]
  \item Severe or worsening symptoms of eosinophilic fasciitis (shulman syndrome) require immediate medical evaluation
\end{itemize}
\section{Progression and Stages}
The condition may progress through different stages of severity over time. Early stages often involve mild or subtle symptoms that may be overlooked. Moderate stages involve more noticeable symptoms affecting daily function. Advanced stages may involve significant impairment and complications.
\section{Complications}
\begin{itemize}[noitemsep]
  \item Disease progression leading to worsening symptoms
  \item Secondary infections or injuries
  \item Side effects of treatment
  \item Reduced quality of life
  \item Emotional and psychological impact
\end{itemize}
\section{Diagnosis}
To make the diagnosis, doctors look for a deep incisional biopsy including skin, subcutaneous tissue, fascia, and muscle, showing fasciitis with a mixed inflammatory infiltrate and often eosinophils, with peripheral eosinophilia ($>$5\% or absolute count $>$500/$\mu$L) present in 70--90\%.

\begin{itemize}[noitemsep]
  \item Comprehensive medical history and physical examination
  \item Laboratory tests to assess organ function
  \item Imaging studies to visualize affected structures
  \item Specialized testing depending on the affected system
  \item Consultation with relevant specialists
\end{itemize}
\section{Prevention}
\begin{itemize}[noitemsep]
  \item Maintain a healthy lifestyle with balanced diet and exercise
  \item Avoid known risk factors
  \item Attend regular medical check-ups for early detection
  \item Follow recommended screening guidelines
\end{itemize}
\section{Treatment Options}
Treatment typically involves with corticosteroids (prednisone 0.5--1 mg/kg/day), with methotrexate, mycophenolate, and azathioprine used as steroid-sparing agents. Symptomatic management and supportive care Enzyme replacement therapy for certain conditions Gene therapy (emerging for select conditions) Dietary modifications (e.g., phenylketonuria) Regular monitoring for associated complications Physical and occupational therapy Genetic counseling for family planning Participation in clinical trials for new therapies

\begin{itemize}[noitemsep]
  \item Medications to manage symptoms and address underlying mechanisms
  \item Lifestyle modifications and self-care measures
  \item Physical therapy and rehabilitation
  \item Surgical intervention when indicated
  \item Regular monitoring and follow-up care
\end{itemize}
\section{Living with It}
\begin{itemize}[noitemsep]
  \item Follow your treatment plan as prescribed
  \item Maintain regular follow-up with your healthcare provider
  \item Make necessary lifestyle adjustments
  \item Seek support from family, friends, or support groups
  \item Stay informed about your condition
\end{itemize}
\section{Outlook}
\begin{itemize}[noitemsep]
  \item The outlook varies from person to person
  \item Most patients improve with corticosteroids, but joint contractures may persist.
\end{itemize}
